% =====================================================================
%  Partial topology identification and cost-aware probing for
%  low-voltage feeders: method and algorithms
%  Paper-style draft v9 (load gate; meter-level labels), 2026-11-26
%
%  Status of the parts (read before quoting anything from this file):
%   * Algorithm 1 is implemented in rnj_wzzt_dr/identification_v8 (self-
%     contained; development only) and smoke-tested on two small
%     development cases; it has NOT been run at scale. Algorithm 2
%     (probing) is PROPOSED and not implemented.
%   * The registered sealed test evaluates the predecessor of Algorithm 1
%     (identification/final.py at commit 71684ff); Section VIII lists the
%     differences. Its result does not carry over to this method, which
%     needs development checks and its own registered test.
%  History: v4 f5299d4, v5 5e36b4c, v6 4881a92, v7 dd8cd4a, v8 52313c3 and
%  b1c5e4d. v9 (user's decision and an external review, 2026-11-26): a load
%  is modelled only when the PCC balance shows one (gate at 1% of the
%  metered power); labels and anchors are meter groupings, the load's place
%  is reported separately; per-meter bounds optional; wide nodes priced by
%  a node MIQP; certificate scopes; re-solve after pricing; a stopping rule
%  with a proof; the load evidence from two optimised models, with bounds.
%
%  Compile with pdfLaTeX (e.g. on Overleaf). No external .bib needed.
% =====================================================================
\documentclass[journal,onecolumn,11pt]{IEEEtran}

\usepackage{amsmath,amssymb}
\usepackage{booktabs}
\usepackage{array}
\usepackage{algorithm}
\usepackage[noend]{algpseudocode}
\usepackage{url}
\usepackage[hidelinks]{hyperref}

\newtheorem{lemma}{Lemma}
\newtheorem{proposition}{Proposition}
\newtheorem{theorem}{Theorem}
\newtheorem{corollary}{Corollary}
\newtheorem{definition}{Definition}
\newtheorem{remark}{Remark}

\newcommand{\tp}{^{\top}}
\newcommand{\one}{\mathbf{1}}
\newcommand{\Lin}{\mathcal{L}}
\newcommand{\Nset}{\mathcal{N}}
\newcommand{\Bset}{\mathcal{B}}
\newcommand{\Sset}{\mathcal{S}}
\newcommand{\Tset}{\mathcal{T}}
\newcommand{\Uset}{\mathcal{U}}
\newcommand{\Aset}{\mathcal{A}}
\newcommand{\Kset}{\mathcal{K}}
\newcommand{\Gset}{\mathcal{G}}
\newcommand{\cross}{\pitchfork}
\algrenewcommand\algorithmicrequire{\textbf{Input:}}
\algrenewcommand\algorithmicensure{\textbf{Output:}}

\begin{document}

\title{Partial Topology Identification and Cost-Aware Probing\\ for Low-Voltage Feeders}
\author{Method draft v9 of 26~November~2026 (authors to be added)}
\maketitle

\begin{abstract}
Smart meters at the customers of a low-voltage feeder rarely determine its whole topology: some lines are fixed by the
data, others are not. We separate three questions. Is there an unmetered load? The balance at the point of common coupling
decides this before any topology is fitted, and a load is modelled only when it is clearly present. Which groupings of
meters does the data decide? A penalised likelihood with a structural prior scores laminar families of meter sets, with the
unmetered load as one more leaf whose path the laminar constraints encode, and one mixed-integer quadratic program (MIQP)
finds the best tree over a pool of candidate lines. A grouping of meters is \emph{decided} when a solve of the same MIQP
without it proves that no tree within a fixed score margin exists, and \emph{open} when such a tree is found; where the
load sits is reported separately. How certain is the computation? Every label carries the solver's interval, and every
certificate its scope. We show that the evidence for a grouping comes only from the voltage couplings across its children,
that it grows linearly with the data for a true grouping and stays bounded for a false one, and that decided groupings
(anchors) reduce the exact search to unions of a few units per anchor. Probing buys demand-response blocks for the open
groupings that feasible blocks can reach, each hour the one block whose expected evidence is worth its price.
\end{abstract}

\begin{IEEEkeywords}
Topology identification, smart meters, partial identification, laminar families, mixed-integer quadratic programming,
unmetered load, active probing, demand response.
\end{IEEEkeywords}

\noindent\textbf{Status of each part.} Algorithm~\ref{alg:identify} is implemented for development and checked on two small
cases only; Algorithm~\ref{alg:probe} is proposed. A registered sealed test evaluates the predecessor of
Algorithm~\ref{alg:identify}; Section~\ref{sec:tests} lists the differences and the runs this version needs.

% ---------------------------------------------------------------------
\section*{Nomenclature}
\begin{tabular}{@{}p{0.27\textwidth}p{0.68\textwidth}@{}}
\toprule
$\Nset=\{1,\dots,n\}$, $t$, $N$ & Metered customers, time slots, rows of data ($N=Tn$)\\
$s$; $\Nset^+$ & Unmetered load, a leaf with known power and no voltage; all leaves ($\Nset\cup\{s\}$ or $\Nset$)\\
$\varsigma$, $\varsigma_0$ & Unmetered share of the metered power (\ref{eq:gate}); gate threshold (1\%)\\
$y_{i,t}=V_{0,t}^2-V_{i,t}^2$ & Squared-voltage drop at meter $i$\\
$C\subseteq\Nset^+$; $z_C$, $h_C$; $C\cross D$ & Support of a line; its meters, whether it carries the load; $C$ and $D$ cross\\
$\pi(C)=C\cap\Nset$; $\Gset(F)$ & Meters of a support; the groupings of a family (Section~\ref{sec:score})\\
$F\supseteq\Bset$, $\Lin(F)$, $k(F)$ & Topology as a laminar family; its lines; their number\\
$\Omega_0$; $\Omega$ & Initial candidate lines (Appendix~\ref{app:pool}); current pool\\
$H(F)$; $H^{\rm f}_M$ & Score (\ref{eq:score}), half-BIC units; flip problem (\ref{eq:flip})\\
$m=\ln100\approx4.6$ & Score margin of labels and anchors\\
$e_M$, $F_M$ & Gap of grouping $M$ (\ref{eq:gap}) and its rival tree\\
$\Aset$; $\Kset_A$, $k_A$; $K$ & Anchors; units of anchor $A$ and their number; largest node searched by unions\\
$\Tset$; $u$, $\Uset_h$, $c(u)$; $\psi$ & Probing target set; block, feasible blocks, price; value of one unit of evidence\\
\bottomrule
\end{tabular}

% ---------------------------------------------------------------------
\section{Model}\label{sec:model}

A feeder has a root, the point of common coupling (PCC), and $n$ metered customers. Meter $i$ reports $p_{i,t}$,
$q_{i,t}$, $V_{i,t}$; the PCC reports $V_{0,t}$ and, where metered, $P^{\rm pcc}_t$, $Q^{\rm pcc}_t$. Every line feeds a set
of leaves, its \emph{support}; the supports of a tree are laminar, and a laminar family containing
$\Bset=\{\{1\},\dots,\{n\},\Nset^+\}$ fixes the tree up to internal nodes without leaves, which voltages cannot see: only the
summed impedance through a node of degree two can be recovered. A topology is a laminar family $F\supseteq\Bset$ over
$\Nset^+$ with lines $\Lin(F)=F\setminus\Bset$. With the linearised DistFlow equations,
\begin{equation}\label{eq:model}
y_t=\sum_{C\in F} z_C\big[r_C\,(z_C\tp p_t+h_C u^P_t)+x_C\,(z_C\tp q_t+h_C u^Q_t)\big]+\varepsilon_{0,t}\one+o+e_t ,
\end{equation}
with $z_C\in\{0,1\}^n$ the meters of $C$, $h_C=1$ if $s\in C$, $0\le r_C,x_C\le\bar\theta$ twice the line's resistance and
reactance in pu (nuisance; $\bar\theta=2$ never binds), $\varepsilon_{0,t}$ the error of the PCC voltage reading, $o$ fixed
meter offsets and $e_t$ meter noise. One unmetered load is allowed. Its power is the PCC balance minus the losses estimated
from the meters' own voltages,
\begin{equation}\label{eq:balance}
u^P_t=P^{\rm pcc}_t-\sum_i p_{i,t}\Big(1+\frac{V_{0,t}-V_{i,t}}{V_{i,t}}\Big),\qquad u^Q_t\ \text{likewise}.
\end{equation}
When the gate of Section~\ref{sec:gate} finds a load, it is one more leaf $s$, with known power and no voltage reading;
otherwise there is no leaf $s$ and $h\equiv0$. The supports that carry the load all contain $s$, so in a laminar family they
are nested: they are the path from the PCC to the load (Proposition~\ref{prop:path}). The admissible places of the load are
below any line, below the trunk (the load shares the trunk), and at the PCC (beside the trunk: the meters' trunk is then the
line $\Nset$ without $s$). The load's own branch $\{s\}$ affects no voltage: its columns are zero, it carries no price, it is
compatible with every support, and in pricing it is a child like any leaf. In matrix form, $y_t=Rp^+_t+Xq^+_t+\dots$ with
$p^+_t=(p_t,u^P_t)$ and $R=\sum_Cr_Cz_Cz_C^{+\top}$, $z^+_C$ the leaves of $C$: $R_{ij}$ is the resistance shared by the
paths of meter $i$ and leaf $j$. The data come from AC power flow; (\ref{eq:model}) is used inside the method only.
\emph{Assumptions stated, not modelled:} the loss estimate (\ref{eq:balance}) is approximate, several unmetered loads with
different profiles are not one aggregate, and errors in the measured $p,q$ enter the regressors.

% ---------------------------------------------------------------------
\section{Is there an unmetered load?}\label{sec:gate}

Before any topology is fitted, the PCC balance decides whether a load is modelled at all. With $u^P$ from
(\ref{eq:balance}),
\begin{equation}\label{eq:gate}
\varsigma=\frac{\sum_tu^P_t}{\sum_t\sum_ip_{i,t}},\qquad\text{a load is modelled if and only if }\varsigma\ge\varsigma_0=1\%.
\end{equation}
This is a prior setting, not a test inside the fit: a load whose mean power is below 1\% of the metered power is treated as
absent. Fitting a load of about zero would be wrong, not merely useless: a line $M$ and the same line carrying the load
would fit equally well, and the load's path would be chosen by the noise of $u$. On the 288 development cases the statistic
separates cleanly (Table~\ref{tab:gate}): without a load it stays within $-0.49\%$ and $0.66\%$; the smallest load is
$1.53\%$. The threshold is fixed from these cases before any test of the method, and the gate's error rates (loads missed,
no-load feeders given a load) are reported by load share and kind. The spread without a load comes from the meters' gain
errors (0.3\%, 0.2\% at the PCC) and the loss estimate, so other meter classes need their own calibration.
\begin{table}[h]
\caption{Unmetered share $\varsigma$ on the development cases (no identifier involved)}\label{tab:gate}
\centering
\begin{tabular}{@{}lrr@{}}
\toprule
Cases & Number & $\varsigma$: min / median / max\\
\midrule
No load & 72 & $-0.49\%$ / $0.03\%$ / $0.66\%$\\
Load of 2\% (three kinds) & 72 & $1.53\%$ / $2.0\%$ / $2.61\%$\\
Load of 5\% & 72 & $4.59\%$ / $5.0\%$ / $5.81\%$\\
Load of 10\% & 72 & $9.51\%$ / $10.2\%$ / $11.68\%$\\
\bottomrule
\end{tabular}
\end{table}

% ---------------------------------------------------------------------
\section{Score and labels}\label{sec:score}

\paragraph*{Score} The nuisance terms $o$ and $\varepsilon_{0,t}\one$ are removed by $M=D\otimes S_\gamma$, with
$D=I_T-\Pi_{[\one]}$ and $S_\gamma=I_n-\gamma\one\one\tp/n$ ($\gamma$ the generalised least squares weight of the shared PCC
error). The score is
\begin{equation}\label{eq:score}
H(F)=\min_{\theta}\frac{\|M(y-\Phi_F\theta)\|^2}{2\hat\sigma^2}+\beta\,k(F),\qquad\beta=\tfrac12p_s\ln N-\lambda,
\end{equation}
with $p_s\in\{1,2\}$ identifiable parameters per line (2 when most meters' $q$ is not a fixed multiple of their $p$ within
the meter class). $D$ and $S_\gamma$ act on different axes, and minimising over $o$ and $\varepsilon_0$ gives
(\ref{eq:score}) exactly (Frisch--Waugh--Lovell). \emph{Free load term.} If the load's term in (\ref{eq:model}) is replaced by
$\rho_iu^P_t+\chi_iu^Q_t$ with free $\rho,\chi\in\mathbb R^n$, minimising over them replaces $D$ by
\begin{equation}\label{eq:free}
D_1=I_T-\Pi_{[\one,\,u^P,\,u^Q]}
\end{equation}
for every family at once. \emph{Noise variance.} $\hat\sigma^2$ is the residual variance of the saturated model (every
$R_{ij}$, $X_{ij}$ free) under $M$, and under (\ref{eq:free}) when a load is modelled, so that it does not depend on the
load's place; its degrees of freedom are the rank of the operator minus the effective (numerical) rank of the saturated
design. This assumes noise that is white across slots and meters after $M$; unequal or correlated meter errors are not
modelled, which is why every threshold is calibrated on data rather than taken from asymptotics.
\begin{proposition}[Structural prior]\label{prop:prior}
Under the Gibbs prior $P(F)\propto\prod_v\rho^{c_v-2}$ over laminar families ($c_v$ the children of internal node $v$),
$-\ln P(F)=\mathrm{const}-\lambda\,k(F)$ with $\lambda=\ln(1/\rho)$.
\end{proposition}
So score differences approximate log posterior odds. We set $\rho=0.141$, the odds of a branching point having more than two
children on 84 real LV feeders ($\lambda=1.96$); the same feeders set the statistics of our case generator, so results with
$\lambda=0$ are reported alongside.

\paragraph*{Groupings} With a load, one grouping of meters can appear as two supports, $M$ and $M\cup\{s\}$: the same
statement about the meters, with the load inside the line or not. Write $\pi(C)=C\cap\Nset$ and call
$\Gset(F)=\{\pi(C):C\in F,\ 2\le|\pi(C)|\le n-1\}$ the \emph{groupings} of $F$: its meter topology. Two topologies with the
same groupings differ only in where the load sits. Labels, anchors and probing concern groupings; the load's place is
reported separately (Section~\ref{sec:id}). Without a load, $\pi$ is the identity and groupings are lines.

\paragraph*{Competitive set and labels} For a pool $\Omega$ and $H^*=\min_FH(F)$,
$\Sset_m=\{F\subseteq\Omega\cup\Bset\ \text{laminar}:\ H(F)\le H^*+m\}$. For a tree $T$ and a grouping $M$, the \emph{flip
problem} asks for the best tree that decides $M$ the other way,
\begin{equation}\label{eq:flip}
H^{\rm f}_M(T)=\min\{H(F):\ F\subseteq\Omega\cup\Bset\ \text{laminar},\ [M\in\Gset(F)]\neq[M\in\Gset(T)]\},
\end{equation}
which for $M\in\Gset(T)$ is the tree MIQP (\ref{eq:master}) with every support of grouping $M$ removed from the pool. Its
\emph{gap} is
\begin{equation}\label{eq:gap}
e_M(T)=H^{\rm f}_M(T)-H(T)\ \in\ \big[\,L_M-H(T),\ U_M-H(T)\,\big],
\end{equation}
with $L_M$ the solver's lower bound and $U_M$ the score of the best flipped tree found, the \emph{rival} $F_M$.
\begin{definition}[Labels]
Relative to $T$, grouping $M$ is \emph{decided} if $L_M>H(T)+m$, \emph{open} if a flipped tree with $H\le H(T)+m$ is found,
and \emph{unsettled} otherwise. The \emph{flip test} solves (\ref{eq:flip}) with the threshold $H(T)+m$ as the solver's
cutoff; it stops as soon as ``decided'' is proved, and otherwise runs to its limit to find the best rival.
\end{definition}
\begin{proposition}[What a label certifies]\label{prop:label}
If $M$ is decided, every tree in $\Sset_m$ agrees with $T$ on whether $M$ is a grouping. If $T$ is a proven optimum and $M$
is open, $\Sset_m$ contains trees that disagree.
\end{proposition}
Labels are statements about the model and the pool, not about the truth. They separate three cases: certified decided,
certified open, and unsettled (computation unresolved); the interval (\ref{eq:gap}) carries the solver's uncertainty into
probing.

\paragraph*{Local flips} Removing the lines of grouping $M$ from $T$, or swapping in a pool line $D$ whose grouping crosses
$M$ (the lines of $T$ that cross $D$ removed), gives flipped trees without solving an MIQP:
\begin{equation}\label{eq:local}
m_M(T)=\min\Big\{H(T\setminus\pi^{-1}(M)),\ \min_{D:\ \pi(D)\cross M}H(T_{|D}\cup\{D\})\Big\}-H(T),\qquad
\ell_D(T)=H(T)-H(T_{|D}\cup\{D\}),
\end{equation}
with $T_{|D}=\{E\in T:\neg(E\cross D)\}$.
\begin{lemma}[Direction of the local bounds]\label{lem:local}
$e_M(T)\le m_M(T)$ for $M\in\Gset(T)$, and $e_{\pi(D)}(T)\le-\ell_D(T)$ for $\pi(D)\notin\Gset(T)$. Local flips can prove a
grouping open, never decided.
\end{lemma}

% ---------------------------------------------------------------------
\section{Evidence for one line}\label{sec:evidence}

Let $C\in\Lin(T)$ have children $A_1,\dots,A_c$ in $T$ ($c\ge2$): the largest supports of $T$ inside $C$ and the leaves of
$C$ outside them. Its \emph{cross-child pairs} are $\times(C)=\{(i,j):\ i\in A_a\cap\Nset,\ j\in A_b,\ a\ne b\}$.
\begin{lemma}[Glue]\label{lem:glue}
With $Z^\times_C$ the indicator of $\times(C)$,
\begin{equation}\label{eq:glue}
r_C\,z_Cz_C^{+\top}=\sum_{a=1}^{c}r_C\,z_{A_a}z_{A_a}^{+\top}+r_C\,Z^\times_C .
\end{equation}
Removing $C$ and adding $r_C,x_C$ to every child reproduces every prediction of $T$ except the entries on $\times(C)$, and no
other line of $T$ can change those entries alone. The evidence for $C$ therefore comes only from its cross-child pairs; a
meter's own cable has none, so a line needs at least two leaves to carry evidence.
\end{lemma}
\begin{lemma}[Removal gap]\label{lem:wald}
For $C\in\Lin(T)$, with the bounds inactive at the fitted parameters,
\begin{equation}\label{eq:wald}
H(T\setminus\{C\})-H(T)=\frac{\hat\theta_C\tp S_C\,\hat\theta_C}{2\hat\sigma^2}-\beta,\qquad
S_C=\Phi^{\times\top}_C\big(I-\Pi_{T\setminus\{C\}}\big)\Phi^\times_C ,
\end{equation}
where $\Phi^\times_C$ are the columns of $C$ in which meter $i$ sees only the flows of the leaves of $C$ outside its own child.
The first term is half the Wald statistic of $C$: the information about $C$ left after every other line.
\end{lemma}
\begin{theorem}[Evidence grows for true lines only]\label{thm:gap}
Let $n$ be fixed and $T\to\infty$. Assume (A1) the data follow (\ref{eq:model}) with tree $T^0$ and noise that is white under
$M$ with variance $\sigma^2$, and $\hat\sigma^2\to\sigma^2$; (A2) every line of $T^0$ has $r_C\ge r_{\min}>0$ and its
parameters lie inside the bounds; (A3) for every $C\in\Lin(T^0)$,
\[
\kappa_C=\liminf_{N\to\infty}\ \frac1{N\sigma^2}\ \min_{F\not\ni C}\big\|(I-\Pi_F)\,M\Phi_{T^0}\theta^0\big\|^2>0,
\]
the minimum over the finitely many laminar families without $C$. Then (a) if $C\in\Lin(T^0)$,
$H(F)-H(T^0)\ge\tfrac12N\kappa_C-O_p(\sqrt N)-O(\ln N)$ uniformly over families $F\not\ni C$, so $e_C(T^0)\to\infty$ and
$m_C(T^0)\to\infty$ in probability, for every pool; (b) if $C\notin T^0$ and $T\supseteq T^0\cup\{C\}$ is laminar,
$H(T\setminus\{C\})-H(T)\le O_p(1)-\beta\to-\infty$.
\end{theorem}
With a load, the same statements hold for groupings, with ``families without $C$'' read as ``families without grouping $M$''
in (A3); the proof is unchanged.
\begin{corollary}[Anchors are true groupings]\label{cor:true}
If $\Lin(T^0)\subseteq\Omega$, then with probability tending to one $T^0$ is the best tree over $\Omega$ and every grouping
of $T^0$ is decided, for any fixed $m$.
\end{corollary}
\begin{remark}
(A3) fails for voltage twins, whose couplings another family reproduces: their gap does not grow and they stay open, which
is the right answer. (A3) is a property of the loads, and probing can create it (Section~\ref{sec:probe}).
\end{remark}

% ---------------------------------------------------------------------
\section{Identification}\label{sec:id}

\paragraph*{Tree MIQP} For a pool $\Omega$ of supports, with a binary $b_C$ per support,
\begin{equation}\label{eq:master}
\begin{aligned}
\min_{b,r,x}\ &\frac{\|M(y-\Phi(r,x))\|^2}{2\hat\sigma^2}+\beta\sum_{C\in\Omega}b_C\\
\text{s.t.}\ \ &0\le r_C\le\bar r_Cb_C,\quad0\le x_C\le\bar x_Cb_C\quad(C\in\Omega),\qquad b_C+b_D\le1\quad(C\cross D),\\
&\textstyle\sum_{C\in\Omega,\,C\subsetneq W}b_C\le|W|-2\quad(W\in\Omega\cup\{\Nset^+\}),\qquad
\sum_{C\in\Omega,\ \pi(C)=A}b_C\ge1\quad(A\in\Aset),\\
&\textstyle\sum_{C\ni i}r_C\le\bar R_i,\quad\sum_{C\ni i}x_C\le\bar X_i\quad(i\in\Nset;\ \text{optional}),
\end{aligned}
\end{equation}
with $\Bset$ always on. The count cut holds for every laminar family (a tree on $|W|$ leaves has at most $|W|-2$ lines
strictly inside $W$). The fourth constraint keeps every anchor's grouping, with the load inside or not. Given a level $\ell$
that some tree reaches (an incumbent's score, or a flip test's threshold), $\bar R_i$ and $\bar X_i$ are the largest
diagonal responses of any coefficients whose residual fits under $\ell$, from one convex quadratically constrained program
per meter with the binaries dropped, and $\bar r_C=\min(\bar\theta,\min_{i\in\pi(C)}\bar R_i)$. These are valid
inequalities: every tree with $H\le\ell$ satisfies them, so they change neither the optimum nor a flip test's verdict when
the solver finishes; they only strengthen the relaxation, and their effect is measured by an ablation. Each solve starts
from a local-search incumbent and keeps the better tree.
\begin{proposition}[The load's path needs no constraints]\label{prop:path}
For a set of meters $C$ write $C^1=C\cup\{s\}$. Over $\Nset^+$ the crossing constraints of (\ref{eq:master}) are exactly:
(i) the crossings of the meter sets; (ii) one place: $b_{C^1}+b_{D^1}\le1$ unless $C\subseteq D$ or $D\subseteq C$; (iii) a
line above a line that carries the load carries it too: $b_{C^1}+b_D\le1$ for meter sets $C\subsetneq D$. The indicator
formulation ($h_C\le b_C$; $h_C+b_D-h_D\le1$ for $C\subsetneq D$; $h_C+h_D\le1$ for $C,D$ not nested; products $r_Ch_C$) is
the special case in which no meter set appears both with and without $s$; the leaf form needs neither its constraints nor its
products, and also lets the load join in the middle of a line.
\end{proposition}

\paragraph*{Node pricing} Let $v$ be an internal node of $T$ with children $A_1,\dots,A_{c_v}$ (supports or leaves).
\begin{lemma}[Node decomposition]\label{lem:node}
A support $C\notin T$ with at least two leaves is compatible with $T$ if and only if $C=\bigcup_{j\in S}A_j$ for a unique
internal node $v$ and a set $S$ of its children with $2\le|S|<c_v$.
\end{lemma}
\begin{equation}\label{eq:price}
g_v(T)=\max_{S:\ 2\le|S|<c_v}H(T)-H\Big(T\cup\Big\{\textstyle\bigcup_{j\in S}A_j\Big\}\Big).
\end{equation}
A node with at most 12 children is priced by enumeration: every union is fitted. A wider node is priced by a node MIQP with
a binary $z_a$ per child; the new line's coefficient enters as $u_{ab}=rz_az_b$, linearised exactly by McCormick inequalities
($u_{ab}\le r$, $u_{ab}\le\bar\theta z_a$, $u_{ab}\le\bar\theta z_b$, $u_{ab}\ge r-\bar\theta(2-z_a-z_b)$), so that meters of
child $a$ see the flow of child $b$; every coefficient of $T$ is refitted. Both give $\max_vg_v$ exactly when they finish. We
distinguish three \emph{certificate scopes}: (a) \emph{complete}: every node was enumerated or its MIQP proved optimal, and no
line pays, so no single line compatible with $T$, among all supports, lowers the score; (b) \emph{restricted}: statements over
the anchored universe below, which covers every combination of new lines inside anchor nodes with at most $K$ units (so pair
moves need no separate search there) but not inside wider nodes; (c) \emph{incomplete}: a node MIQP stopped at its time
limit. Neither certificate shows that no line is missing: a wrong tree may need lines removed, and a node with two children
has nothing to price. We report the scope, the number of nodes priced each way, and the share of internal nodes with at
least three children.

\paragraph*{Anchors}
\begin{definition}[Anchors and units]
An \emph{anchor} is a decided grouping of $T$. For a laminar set of anchors $\Aset$ (sets of meters) and each
$A\in\Aset\cup\{\Nset\}$, the \emph{units} $\Kset_A$ are the largest anchors strictly inside $A$ and the meters of $A$ outside
them; $k_A=|\Kset_A|$. For $K\ge2$ the \emph{anchored universe} is the set of meter lines
\begin{equation}\label{eq:universe}
\Omega_K(\Aset)=\Aset\ \cup\bigcup_{A:\,k_A\le K}\Big\{\textstyle\bigcup_{B\in S}B:\ S\subseteq\Kset_A,\ 2\le|S|\le k_A-1\Big\}
\ \cup\ \Omega_{>K}(\Aset),
\end{equation}
with $\Omega_{>K}(\Aset)$ the groupings of the pool compatible with $\Aset$ whose smallest enclosing member of
$\Aset\cup\{\Nset\}$ has more than $K$ units. With a load, $\Omega^+_K(\Aset)$ is every line of $\Omega_K(\Aset)$ with and
without $s$, the supports $\{i,s\}$, and $\Nset$; without a load, $\Omega^+_K=\Omega_K$.
\end{definition}
\begin{proposition}[Anchored universe]\label{prop:anchor}
Let $F$ be a laminar family over $\Nset^+$ with $\Aset\subseteq\Gset(F)$. Every support $C\in\Lin(F)$ has $\pi(C)$ equal to an
anchor, to a union of 2 to $k_A-1$ units of one $A\in\Aset\cup\{\Nset\}$, to a single meter, or to $\Nset$. Hence, if every
$k_A\le K$, $\Lin(F)\subseteq\Omega^+_K(\Aset)$, and (\ref{eq:master}) over $\Omega^+_K(\Aset)$ with the anchors' groupings
required returns the best family that keeps every anchor, among all supports and every admissible place of the load. Its
union binaries number $\sum_A(2^{k_A}-k_A-2)$ (twice that with a load) instead of $2^n-n-2$.
\end{proposition}
Proposition~\ref{prop:anchor} turns ``the anchors are right'' into ``the search is exact''. The flip tests check that premise
over the pool, and by Theorem~\ref{thm:gap} and Corollary~\ref{cor:true} it holds with probability tending to one. The
anchors put the exact search where the data have not decided; a large anchor is refined, not replaced, since its own anchors
become its units. A node with more than $K$ units keeps the pool's lines inside it and is priced by (\ref{eq:price}).

\paragraph*{Pricing, re-solve and stopping} Pricing adds a paying line to the tree and to the pool (the universe of later
solves); when it has added any, (\ref{eq:master}) is solved again over the universe and the tree, and the better tree is
kept.
\begin{lemma}[Re-solve and stopping]\label{lem:stop}
(i) The score never rises in a round. (ii) If a round leaves the anchors unchanged, its pricing adds no line and finds none
that pays, and its MIQP ends optimal, then a further round cannot lower the score.
\end{lemma}

\paragraph*{Load evidence} Whether a load is modelled is decided by the gate (\ref{eq:gate}), not by the fit. The fit
reports, as a diagnostic, the evidence that the voltages show the load: with $H_0$ the score of the same data and
$\hat\sigma^2$ with $u\equiv0$, minimised over meter topologies, and $H$ minimised with the load,
\begin{equation}\label{eq:evidence}
D^*=\min H_0-\min H,\qquad L_0-U_1\ \le\ D^*\ \le\ U_0-L_1,
\end{equation}
when the two minima are bracketed by $[L_0,U_0]$ and $[L_1,U_1]$ (solver bounds and incumbents).
\begin{proposition}[Evidence from two optimised models]\label{prop:evidence}
Evaluating $H_0$ at the meter topology of the load-present tree, $\pi(\hat T)$, can only overstate $D^*$ for the same
$\hat T$, since $\min H_0\le H_0(\pi(\hat T))$. The bounds of (\ref{eq:evidence}) hold for any brackets; we report the
estimate $U_0-U_1$ and the lower end $L_0-U_1$, both over the pools searched.
\end{proposition}
The places of the load are the admissible places in $\hat T$'s meter topology whose score is within $m$ of the best.

\begin{algorithm}[t]
\caption{\textsc{Identify}: gate, tree, anchors and labels}\label{alg:identify}
\begin{algorithmic}[1]
\Require data; $\varsigma_0$, $\lambda$, $m$, $K$, $R$
\Ensure the gate's decision; tree $\hat T$ (groupings, and with a load its path); certificate and scope; anchors; pool;
labels with gap intervals and rivals; with a load, its places and evidence interval
\State \textbf{if} the PCC power is metered: $\varsigma$ by (\ref{eq:gate}); a load is modelled iff $\varsigma\ge\varsigma_0$ \Comment{gate}
\State $M$, $\hat\sigma^2$, $\beta$; $\hat R$ by (\ref{eq:rhat}) \Comment{set-up}
\State $\Omega\gets\Omega_0$ and nine starting trees \Comment{Appendix~\ref{app:pool}}
\State $T\gets$ (\ref{eq:master}) over $\Omega$, from a local-search incumbent
\For{$r=0,1,\dots,R$}
  \State label the groupings of $T$: local flips (\ref{eq:local}), then flip tests (\ref{eq:flip}); $\Aset\gets$ the decided ones
  \State \textbf{if} the stopping rule of Lemma~\ref{lem:stop} holds, or $r=R$: \textbf{break}
  \State $T\gets$ (\ref{eq:master}) over $\Omega^+_K(\Aset)\cup\Lin(T)$, anchors required; $\Omega\gets\Omega\cup\Omega^+_K(\Aset)$
  \State pricing (\ref{eq:price}): add paying lines to $T$ and $\Omega$; if any, (\ref{eq:master}) again; keep the better tree
\EndFor
\State $\hat T\gets T$; groupings outside $\hat T$ whose swap (\ref{eq:local}) is within $m$ are open
\State with a load: its path, its places within $m$, its evidence (\ref{eq:evidence})
\end{algorithmic}
\end{algorithm}

\paragraph*{Explanation} Line 1 decides whether there is a load; below the gate, the rest runs on the meters alone. Lines 2--4
build a first tree from cheap candidate lines. Lines 5--9 are the core: decided groupings become anchors; the anchored
universe offers every line compatible with them at the nodes that are small enough, which is exact by
Proposition~\ref{prop:anchor}; pricing searches the wider nodes and checks the result against any single line, proposed or
not; the loop stops by Lemma~\ref{lem:stop} or after $R$ searches, and the output says which. Lines 10--11 report what remains
open and, with a load, where it sits and how strongly the voltages show it.

% ---------------------------------------------------------------------
\section{Probing}\label{sec:probe}

The operator can buy signed active-power blocks at chosen meters from an aggregator of household batteries and EV chargers,
each held for one market hour and sealed one hour ahead. Probing works on the pool, the leader and the open groupings; it
never lists whole trees. A block that changes the flows across a grouping's children raises its information $S_C$
(Lemma~\ref{lem:wald}): this is what probing buys.

\paragraph*{Values used in probing} For an open grouping $M$ with rival $F_M$, at the gate of hour $h$:
\begin{itemize}
\item the gap interval (\ref{eq:gap}) and the \emph{deficit} $\delta_M=m-(U_M-H(\hat L))$, from the upper end, so that the
deficit is never overstated; unsettled groupings are not targets (probing cannot repair a solver's time limit) and are tested
again at each gate with the new data;
\item the fitted parameters $\hat\theta$ of the leader and the rival, and their covariances $P=\hat\sigma^2(\Phi\tp\Phi)^{-1}$
on the coordinates not at a bound;
\item the predicted voltage drops of hour $h+1$ under each tree, from the day-ahead forecast $\tilde x$ of the meters' load plus
the block $u$. Since (\ref{eq:model}) is linear in the load, their difference is $d_M(u)=d_M(0)+G_Mu$, and
\begin{equation}\label{eq:sep}
q_M(u)=d_M(u)\tp\Sigma_M^{-1}d_M(u),\qquad \Sigma_M=\hat\sigma^2I+\tfrac12\Phi_{\hat L}P_{\hat L}\Phi_{\hat L}\tp+\tfrac12\Phi_{F_M}P_{F_M}\Phi_{F_M}\tp,
\end{equation}
with $\Sigma_M$ evaluated at $\tilde x$;
\item the feasible blocks $\Uset_{h+1}$: battery and EV headroom, the voltage band with a chance margin for the leader
\emph{and every rival}, the transformer rating, and the price caps; the price
$c(u)=c^\uparrow\sum_i\max(u_i,0)+c^\downarrow\sum_i\max(-u_i,0)$.
\end{itemize}
\begin{lemma}[Expected evidence]\label{lem:evidence}
If the next hour's data are Gaussian with covariance $\Sigma_M$ and mean given by the leader (or the rival), the expected
change of $H(F_M)-H(\hat L)$ over that hour is $+\tfrac12q_M(u)$ (or $-\tfrac12q_M(u)$).
\end{lemma}
\begin{lemma}[Reach]\label{lem:reach}
If $G_M=0$, no block changes $q_M$. (i) If the two trees differ only in grouping $M$, $G_M=0$ exactly when no meter with a
controllable battery or EV lies below $M$. (ii) Trees that differ only in where the load sits have $G=0$: the load's place
is decided by the natural variation of the load, never by probing.
\end{lemma}

\paragraph*{Target set and objective} At the start of the day, $\Tset$ is the set of open groupings with $G_M\ne0$; it stays
fixed during the day. The block for hour $h+1$ is
\begin{equation}\label{eq:block}
u^\ast=\arg\max_{u\in\Uset_{h+1}}\ J_h(u),\qquad
J_h(u)=\psi\sum_{M\in\Tset\ \text{open}}\Big[\min\big(\delta_M,\tfrac12q_M(u)\big)-\min\big(\delta_M,\tfrac12q_M(0)\big)\Big]-c(u),
\end{equation}
bought only if $J_h(u^\ast)>0$: the expected evidence beyond the natural load variation, counted up to what each grouping
still needs, against the price. \emph{Choice of $\psi$:} on the development cases, Algorithm~\ref{alg:probe} runs for every
$\psi$ on a grid; with mean cost and coverage scaled to $[0,1]$, the default is the grid point farthest above the straight line
from the first point to the last (the knee). Grid and default are fixed before any probing test; the whole curve is reported.
\emph{Greedy, in plain words:} each hour the rule buys the one block worth most for the next hour alone; it chooses again
when the next hour's data arrive. \emph{Stopping and trust gate:} the day ends when every target is decided at two
consecutive gates; blocks are bought only when the leader fits as well as the meter noise allows,
$s^2_{\hat L}/\sigma^2_y\le\kappa_A$.

\begin{algorithm}[t]
\caption{\textsc{Probe}: one day of cost-aware probing (proposed)}\label{alg:probe}
\begin{algorithmic}[1]
\Require the output of Algorithm~\ref{alg:identify}; history; contract; $\psi$, $m$, $\kappa_A$
\Ensure the leader, with every grouping decided, open or unsettled
\State $\Tset\gets$ open groupings with $G_M\ne0$ \Comment{fixed for the day; Lemma~\ref{lem:reach}}
\For{$h=0,1,\dots,23$} \Comment{gate at the start of hour $h$; block for $h+1$ sealed now}
  \State $\hat L\gets$ Algorithm~\ref{alg:identify}, lines 5--9, on all data so far
  \State relabel the groupings of $\Tset$ that are still open: their gap intervals, deficits and rivals
  \State \textbf{if} every grouping of $\Tset$ is decided at two consecutive gates: \Return $\hat L$ with its labels
  \State $u_{h+1}\gets u^\ast$ of (\ref{eq:block}) if $s^2_{\hat L}/\sigma^2_y\le\kappa_A$ and $J_h(u^\ast)>0$, else $0$
  \State hour $h$ runs with $u_h$; add its readings to the data
\EndFor
\State \Return $\hat L$ with its labels \Comment{some targets still open}
\end{algorithmic}
\end{algorithm}

% ---------------------------------------------------------------------
\section{How the outputs are valued}\label{sec:value}

\paragraph*{What the output says} Three parts, kept apart. \emph{The load:} the gate's statistic and decision; with a load,
its path, its places within $m$, and the evidence interval. \emph{The structure:} the groupings of $\hat T$, the anchors, and
for every grouping its label, gap interval and rival. \emph{The computation:} the certificate's scope, the nodes priced by
enumeration and by MIQP, the share of unsettled labels, the solver status of every number, and why the loop stopped.

\paragraph*{How it is scored against the truth} The gate: loads missed and no-load feeders given a load, by share and kind.
The structure: precision, recall and F1 over groupings; recall near the root; exact trees; false anchors; label coverage,
false decisions, the share of open groupings that are true, and the share unsettled under each time limit. The load: whether
its true place is within $m$ of the best. Probing: the share of targets decided by the end of the day, false decisions
against passive operation, cost, energy and time; over $\psi$, as a cost--coverage curve. Coverage and false decisions are
always reported together.

% ---------------------------------------------------------------------
\section{Contributions, limits and tests}\label{sec:tests}

\paragraph*{Contributions} (i) A separation of three questions: whether there is an unmetered load (a gate on the PCC
balance), which meter groupings the data decide (labels and anchors), and how certain the computation is (intervals and
certificate scopes), all feeding probing. (ii) One MIQP over laminar families in which the load is a leaf without voltage,
so its path costs no constraint and no product (Proposition~\ref{prop:path}). (iii) The evidence for one line: only the
couplings across its children count, and it grows linearly with the data for a true line and stays bounded for a false one
(Lemmas~\ref{lem:glue}--\ref{lem:wald}, Theorem~\ref{thm:gap}). (iv) Anchors, certified by flipped solves of the same MIQP,
that reduce the exact search to unions of a few units per anchor (Proposition~\ref{prop:anchor}). (v) Probing targeted at
the open groupings that the available batteries and EVs can reach.

\paragraph*{Limits} The gate is a prior: loads under 1\% of the metered power are ignored, and its threshold is calibrated
for one meter class. Theorem~\ref{thm:gap} is asymptotic and assumes the model. Labels and anchors are relative to the pool;
certificates have the scopes above; the evidence interval holds over the pools searched; thresholds are calibrated on data,
since the topology is discrete and, without a load, the load's place is not identified, so standard likelihood-ratio
asymptotics do not apply. The probing objective is a surrogate with fixed parameters.

\paragraph*{Relation to the registered test} The registered sealed test evaluates the predecessor (final.py at 71684ff). It
has no gate (with the PCC power it always models a load), removes the load with (\ref{eq:free}) during the search and places
it at the end, takes blocks from the top-level lines of the first tree without flip tests, searches the unions of at most 9
such blocks once, prices nodes with an MIQP over the meters only, has no labels, and computes the evidence at
$\pi(\hat T)$. Its result will be reported as it is; it does not carry over to Algorithm~\ref{alg:identify}.

\paragraph*{Development checks before a registered test} (V1) Structure: do meter-level labels separate grouping
uncertainty from the load's place? (V2) The gate and the load: error rates by share and kind; the evidence at $\pi(\hat T)$
against the two-model evidence, on no-load feeders in particular. (V3) Scale: feeders of growing size and branching; time
per case. (V4) Pricing: how often nodes are wide, how often a node MIQP stops unproved, and the effect on F1. (V5) Formulation:
with and without the per-meter bounds (root gap, nodes, time, final gap, share of unsettled flip tests at fixed time). (V6)
Certification: decided, open and unsettled shares under several time limits. (V7) Probing: whether better gap intervals give
better probing. (V8) Calibration: whether decided groupings are wrong at a rate consistent with the margin. Results are
reported as recovery and certification, not only as scores.

% ---------------------------------------------------------------------
\appendices
\section{Proofs}

\begin{IEEEproof}[Proof of Proposition~\ref{prop:prior}]
$F$ is a rooted tree with $n^+=|\Nset^+|$ leaves and $k(F)+1$ internal nodes, each with at least two children; every node but
the root has one parent, so $\sum_vc_v=n^++k(F)$, $\sum_v(c_v-2)=n^+-k(F)-2$ and $P(F)\propto\rho^{\,n^+-2}\rho^{-k(F)}$.
\end{IEEEproof}

\begin{IEEEproof}[Proof of Proposition~\ref{prop:label}]
If $L_M>H(T)+m$, every flipped tree $F$ has $H(F)\ge L_M>H(T)+m\ge H^*+m$, so $F\notin\Sset_m$: every tree of $\Sset_m$ decides
$M$ as $T$ does. If $H(T)=H^*$ and a flipped tree has $H\le H(T)+m$, it lies in $\Sset_m$.
\end{IEEEproof}

\begin{IEEEproof}[Proof of Lemma~\ref{lem:local}]
The trees in (\ref{eq:local}) are laminar. Removing every line of grouping $M$ leaves no support with $\pi=M$; swapping in $D$
with $\pi(D)\cross M$ removes every line of $T$ that crosses $D$, among them all lines of grouping $M$ (their meters cross
$\pi(D)$), and adds no support of grouping $M$. Both decide $M$ the other way, so their scores bound $H^{\rm f}_M$ from above;
likewise the swap of $D$ for $\pi(D)\notin\Gset(T)$. An upper bound on a gap below $m$ shows a flipped tree within $m$; it
cannot show that none exists.
\end{IEEEproof}

\begin{IEEEproof}[Proof of Lemma~\ref{lem:glue}]
The children partition $C$, so $z_Cz_C^{+\top}=\sum_{a,b}z_{A_a}z_{A_b}^{+\top}$; the terms $a=b$ give the sum in
(\ref{eq:glue}) and the terms $a\ne b$ give $Z^\times_C$. Every child is still a support of $T\setminus\{C\}$, so adding $r_C$
to its parameter reproduces the sum exactly, and likewise for $x_C$. In $T\setminus\{C\}$ the supports that contain a pair of
$\times(C)$ are those strictly above $C$; each adds the same value to every entry of $C\times C$, so none can change the
cross-child entries alone. A meter's cable $\{i\}$ has only the entry $(i,i)$, which lies in no $\times(C)$.
\end{IEEEproof}

\begin{IEEEproof}[Proof of Lemma~\ref{lem:wald}]
By the Frisch--Waugh--Lovell theorem, $\mathrm{RSS}(T\setminus\{C\})-\mathrm{RSS}(T)=\hat\theta_C\tp\Phi_C\tp(I-\Pi_{T\setminus
\{C\}})\Phi_C\hat\theta_C$. By (\ref{eq:glue}), $\Phi_C=\sum_a\Phi_{A_a}+\Phi^\times_C$ with every $\Phi_{A_a}$ in the column
space of $T\setminus\{C\}$. Divide by $2\hat\sigma^2$ and subtract the price $\beta$ of one line.
\end{IEEEproof}

\begin{IEEEproof}[Proof of Theorem~\ref{thm:gap}]
Let $\mu=M\Phi_{T^0}\theta^0$, so the data under $M$ are $\mu+\tilde\varepsilon$ with $\tilde\varepsilon$ white (A1). By
(A2), with probability tending to one the bounded fits of $T^0$ and of every family containing it are unbounded fits, and a
bounded fit never has a smaller residual than the unbounded one. (a) For $F\not\ni C$,
$\mathrm{RSS}(F)-\mathrm{RSS}(T^0)\ge\|(I-\Pi_F)\mu\|^2+2\langle(I-\Pi_F)\mu,\tilde\varepsilon\rangle+\|\Pi_{T^0}\tilde\varepsilon
\|^2-\|\Pi_F\tilde\varepsilon\|^2$. The first term is at least $N\sigma^2\kappa_C$ for large $N$; the second has mean zero and
variance $O(N)$; the last is $O_p(1)$. The penalties differ by at most $\beta n^+=O(\ln N)$. There are finitely many $F$. Every
flipped family of $C$ omits $C$, and $m_C\ge e_C$ by Lemma~\ref{lem:local}. (b) $T$ and $T\setminus\{C\}$ contain $T^0$, so
$\mu$ lies in both column spaces and the difference of residuals is $\sigma^2$ times a $\chi^2$ with $p_s$ degrees of freedom
(the bounded one is no larger), while $\beta\to\infty$.
\end{IEEEproof}

\begin{IEEEproof}[Proof of Corollary~\ref{cor:true}]
A family $F\ne T^0$ either omits a line of $T^0$, and then $H(F)-H(T^0)\to\infty$ by (a), or strictly contains $T^0$, and then
its residual is lower by $O_p(1)$ while its penalty is higher by at least $\beta\to\infty$. With finitely many $F$, $T^0$ is
the unique best tree with probability tending to one, and every flipped family of a grouping of $T^0$ scores above
$H(T^0)+m$, which a solve to optimality proves.
\end{IEEEproof}

\begin{IEEEproof}[Proof of Proposition~\ref{prop:path}]
$C^1$ and $D^1$ both contain $s$, so they are nested or cross, and nested exactly when $C\subseteq D$ or $D\subseteq C$. For
$C^1$ and $D$ with $s\notin D$: $C^1\subseteq D$ is impossible, $D\subseteq C^1$ holds exactly when $D\subseteq C$, and they
are disjoint exactly when $C\cap D=\emptyset$; so they cross exactly when $C\subsetneq D$ or $C\cross D$. With $b^0_C=b_C-h_C$
and $b^1_C=h_C$, the indicator formulation's constraints become (ii) and (iii), and $r_Ch_C$ becomes the coefficient of $C^1$.
The leaf form also admits $C$ and $C^1$ together.
\end{IEEEproof}

\begin{IEEEproof}[Proof of Lemma~\ref{lem:node}]
If $C=\bigcup_{j\in S}A_j$ with $2\le|S|<c_v$, $C$ is nested with or disjoint from every support of $T$ and equals none of
them. Conversely, let $v$ be the smallest support of $T$ containing a compatible $C\notin T$; no child contains $C$, so every
child is inside $C$ or disjoint from it, the children partition $v$, and $C=\bigcup_{j\in S}A_j$ with $2\le|S|<c_v$.
\end{IEEEproof}

\begin{IEEEproof}[Proof of Proposition~\ref{prop:anchor}]
If two supports are nested or disjoint, so are their meter sets: removing $s$ keeps inclusions and disjointness. Every anchor
$A$ is $\pi$ of a support of $F$, so for every $C\in\Lin(F)$ the meter set $\pi(C)$ is compatible with every anchor. If
$2\le|\pi(C)|\le n-1$, let $A$ be the smallest member of $\Aset\cup\{\Nset\}$ containing $\pi(C)$. A unit of $A$ that is an
anchor $B$ satisfies $B\subseteq\pi(C)$ or $B\cap\pi(C)=\emptyset$, since $\pi(C)\subsetneq B$ would contradict the choice of
$A$; a unit that is a meter is inside $\pi(C)$ or not. So $\pi(C)$ is a union of units of $A$; one unit makes it an anchor,
all units make it $A$, which is then an anchor. Otherwise $|\pi(C)|\le1$, where $C=\{i,s\}$, or $\pi(C)=\Nset$, where $C=\Nset$.
Since $C$ is $\pi(C)$ or $\pi(C)\cup\{s\}$, $C\in\Omega^+_K(\Aset)$ when every $k_A\le K$. The required groups keep every
anchor, so the MIQP's feasible set is exactly the families over $\Omega^+_K(\Aset)$ that keep the anchors, which contains every
such family over all supports.
\end{IEEEproof}

\begin{IEEEproof}[Proof of Lemma~\ref{lem:stop}]
(i) The re-solve's universe contains the tree's lines, and the better of the two trees is kept. (ii) The next round would
search $\Omega^+_K(\Aset)\cup\Lin(T)$ with the same anchors: the unions are the same, and the pool has grown only by the last
universe and the tree's lines, which were all in the last solve, since pricing added nothing. So the next universe is
contained in the last one and contains $T$, which is optimal over the last one; no tree of the next universe scores lower.
\end{IEEEproof}

\begin{IEEEproof}[Proof of Proposition~\ref{prop:evidence}]
$\min H_0\le H_0(\pi(\hat T))$ because $\pi(\hat T)$ is one meter topology. For the interval, $\min H_0\ge L_0$ and
$\min H\le U_1$ give $D^*\ge L_0-U_1$, and $\min H_0\le U_0$ and $\min H\ge L_1$ give $D^*\le U_0-L_1$.
\end{IEEEproof}

\begin{IEEEproof}[Proof of Lemma~\ref{lem:evidence}]
$H(F_M)-H(\hat L)$ changes by the log-likelihood ratio of the new data, the penalties being unchanged. For two Gaussians with
common covariance its expectation under the leader is the Kullback--Leibler divergence $\tfrac12d\tp\Sigma^{-1}d$, and minus
that under the rival.
\end{IEEEproof}

\begin{IEEEproof}[Proof of Lemma~\ref{lem:reach}]
$d_M(u)=d_M(0)+G_Mu$, so $G_M=0$ gives $q_M(u)=q_M(0)$. (i) With the same other parameters, the block's part of the
difference is $r_Mz_M(z_M\tp B)u$, $B$ selecting the meters with a controllable device; it vanishes exactly when $z_M\tp B=0$.
(ii) Trees that differ only in whether lines carry the load predict drops that differ by terms in $u^P$ and $u^Q$ only.
\end{IEEEproof}

\section{Initial candidate lines used in this work}\label{app:pool}

Any tree builders can supply $\Omega_0$; this work uses nine. (1) A sensitivity matrix without tree structure, fitted to the
active-power part of (\ref{eq:model}) by a convex QP solved with Gurobi, with the free load term (\ref{eq:free}) when a load
is modelled ($D_1=D$ otherwise):
\begin{equation}\label{eq:rhat}
\hat R=\arg\min_R\ \big\|(D_1\otimes S_\gamma)\,\mathrm{vec}(Y-PR)\big\|^2\quad\text{s.t.}\quad R=R\tp\ge0,\quad
R_{ij}\le\min(R_{ii},R_{jj}).
\end{equation}
(2) Nine hierarchies over the meters: recursive neighbour joining, neighbour joining and recursive grouping on $\hat R$;
UPGMA on the path distance $d_{ij}=\hat R_{ii}+\hat R_{jj}-2\hat R_{ij}$ and Ward linkage on $\sqrt d$; average-linkage
clusterings on the correlation distance between meters of the rows of $\hat R$, the drops $y$, the drops
$(D_1\otimes S_\gamma)y$, and the residuals after fitting the service cables alone. (3) With a load, it is placed in each of
the nine trees at its best admissible place (below a line, below the trunk, or at the PCC); these trees start the first
solve. (4) $\Omega_0$ is every cluster with 2 to $n-1$ meters in any hierarchy, and with a load each such cluster with and
without $s$, the supports $\{i,s\}$ and $\Nset$, so $|\Omega_0|\le18(n-2)+n+1$.

\section{Settings}

\begin{table}[h]
\caption{Settings of Algorithm~\ref{alg:identify} (development; values from the registered predecessor where it has them)
and of probing}\label{tab:settings}
\centering
\begin{tabular}{@{}p{0.36\textwidth}p{0.56\textwidth}@{}}
\toprule
Setting & Value\\
\midrule
Gate $\varsigma_0$ & 1\% of the metered mean power (prior setting; Table~\ref{tab:gate})\\
$\lambda$; $m$; $\bar\theta$ & 1.96; 4.6 (half-BIC units); 2 pu\\
$K$; anchored searches $R$ & 9 units; at most 2, then the stopping rule of Lemma~\ref{lem:stop} or the limit\\
Tree MIQP & 120 s up to 35 meters (warm start from the incumbent), 240 s beyond (cold); per-meter bounds on, at a quarter
of the limit at most (off in the ablation)\\
Flip test & cutoff $H(T)+m$, local flips first; limit at least that of the tree MIQP; the share at the limit is reported\\
Local search & from the empty tree and the builder trees; 20 s per start (40 s beyond 35 meters); polish 20 s\\
Pricing & enumeration up to 12 children; node MIQP beyond, 60 s; at most 10 additions per round\\
Load evidence & the no-load model solved over the meter groupings of the pool, same limits\\
Solver & Gurobi, one thread per case\\
\midrule
Probing & $c^\downarrow=4$, $c^\uparrow=1.2$ CNY/kWh; caps 20 CNY/h, 100 CNY/day; band $[0.96,1.035]$, $z=2.33$;
$\kappa_A=1.2$; $\psi$ at the knee of the development cost--coverage curve\\
\bottomrule
\end{tabular}
\end{table}

\end{document}
